\documentclass[aspectratio=169]{beamer}

\usepackage[utf8]{inputenc}
\usepackage[T1]{fontenc}
\usepackage[activeacute]{babel}
\usepackage{amsmath, amssymb}
\usepackage{booktabs}
\usepackage{array}
\usepackage{xcolor}
\usepackage{tcolorbox}

\usetheme{Madrid}
\usecolortheme[RGB={30,60,120}]{structure}

\setbeamertemplate{navigation symbols}{}
\setbeamertemplate{footline}{
  \leavevmode
  \hbox{
    \begin{beamercolorbox}[wd=.45\paperwidth,ht=2.5ex,dp=1ex,left,leftskip=1em]{author in head/foot}
      \usebeamerfont{author in head/foot}\insertshortauthor
    \end{beamercolorbox}
    \begin{beamercolorbox}[wd=.35\paperwidth,ht=2.5ex,dp=1ex,center]{title in head/foot}
      \usebeamerfont{title in head/foot}MAD1 -- Unidad 5
    \end{beamercolorbox}
    \begin{beamercolorbox}[wd=.2\paperwidth,ht=2.5ex,dp=1ex,right,rightskip=1em]{date in head/foot}
      \usebeamerfont{date in head/foot}\insertframenumber{} / \inserttotalframenumber
    \end{beamercolorbox}
  }
}

\definecolor{unsazul}{RGB}{30,60,120}
\definecolor{unsazulclaro}{RGB}{220,230,245}
\definecolor{grisclaro}{RGB}{245,245,245}
\definecolor{verdeapp}{RGB}{20,120,60}

\tcbuselibrary{skins, breakable}

\newtcolorbox{aplicacion}{
  colback=verdeapp!8,
  colframe=verdeapp,
  fonttitle=\bfseries\small,
  title={Aplicaciones reales},
  left=4pt, right=4pt, top=4pt, bottom=4pt
}

\newtcolorbox{concepto}{
  colback=unsazulclaro,
  colframe=unsazul,
  fonttitle=\bfseries\small,
  title={Idea clave},
  left=4pt, right=4pt, top=4pt, bottom=4pt
}

\newtcolorbox{atencion}{
  colback=orange!10,
  colframe=orange!80!black,
  fonttitle=\bfseries\small,
  title={Atenci\'on},
  left=4pt, right=4pt, top=4pt, bottom=4pt
}

\title[Unidad 5 -- Clustering]{Unidad 5: Aprendizaje No Supervisado}
\subtitle{Bloque 1 -- Clustering}
\author[MAD1 -- UNS]{Métodos de Análisis de Datos I}
\institute{Departamento de Matemática -- Universidad Nacional del Sur}
\date{}

%---------------------------------------------------------------
\begin{document}
%---------------------------------------------------------------

\begin{frame}
\titlepage
\end{frame}

%---------------------------------------------------------------
\begin{frame}{¿Qué es el clustering?}
\begin{columns}[T]
\column{0.55\textwidth}
\textbf{Problema:} dado un conjunto de observaciones $\{\mathbf{x}_1,\ldots,\mathbf{x}_n\}$ \emph{sin etiquetas}, encontrar grupos homogéneos.

\medskip
\begin{concepto}
Los puntos dentro de un grupo deben ser \textbf{similares entre sí} y \textbf{distintos} de los de otros grupos. Nadie nos dice cuáles son los grupos: emergen de los datos.
\end{concepto}

\medskip
Diferencia con clasificación:
\begin{itemize}
  \item Clasificación: tenemos $y_i$, aprendemos $f(\mathbf{x}) \to y$.
  \item Clustering: no hay $y_i$; la estructura la descubrimos nosotros.
\end{itemize}

\column{0.42\textwidth}
\textbf{Comparativa de métodos:}
\medskip

{\small
\begin{tabular}{lcc}
\toprule
\textbf{Método} & \textbf{Requiere $k$} & \textbf{Outliers} \\
\midrule
K-means    & Sí  & Sensible \\
Jerárquico & No* & Sensible \\
DBSCAN     & No  & Detecta \\
GMM        & Sí  & Moderado \\
\bottomrule
\multicolumn{3}{l}{\tiny *elige $k$ al cortar el dendrograma}
\end{tabular}
}

\medskip
\begin{aplicacion}
Segmentación de clientes, agrupamiento de genes, detección de tópicos, compresión de imágenes.
\end{aplicacion}
\end{columns}
\end{frame}

%---------------------------------------------------------------
\begin{frame}{Métricas de distancia}
Todo algoritmo de clustering necesita medir ``similitud''. La elección importa.

\medskip
\begin{columns}[T]
\column{0.48\textwidth}
\textbf{Distancia Euclidiana}
\[
d_E(\mathbf{x},\mathbf{z}) = \sqrt{\sum_{j=1}^{p}(x_j - z_j)^2}
\]
La más usada. Sensible a escala.

\medskip
\textbf{Distancia Manhattan}
\[
d_M(\mathbf{x},\mathbf{z}) = \sum_{j=1}^{p}|x_j - z_j|
\]
Más robusta a outliers.

\column{0.48\textwidth}
\textbf{Distancia Coseno}
\[
d_{\cos}(\mathbf{x},\mathbf{z}) = 1 - \frac{\mathbf{x} \cdot \mathbf{z}}{\|\mathbf{x}\|\|\mathbf{z}\|}
\]
Ideal para texto: importa la dirección, no la magnitud.

\medskip
\begin{atencion}
\textbf{Siempre escalar antes de clusterizar.} Si una variable es salario en pesos y otra es edad en años, la primera domina la distancia sin razón conceptual.
\end{atencion}
\end{columns}
\end{frame}

%---------------------------------------------------------------
\begin{frame}{K-means: algoritmo}
\begin{columns}[T]
\column{0.52\textwidth}
\textbf{Objetivo:} minimizar la inercia (varianza intra-cluster):
\[
W = \sum_{l=1}^{k}\sum_{\mathbf{x}_i \in C_l} \|\mathbf{x}_i - \boldsymbol{\mu}_l\|^2
\]

\medskip
\textbf{Algoritmo:}
\begin{enumerate}
  \item Inicializar $k$ centroides.
  \item \textbf{Asignar} cada punto al centroide más cercano.
  \item \textbf{Actualizar} cada centroide como la media de su cluster.
  \item Repetir hasta convergencia.
\end{enumerate}

\medskip
\begin{concepto}
K-means converge siempre, pero puede quedar en un mínimo local. Se recomienda múltiples inicializaciones y quedarse con la de menor inercia.
\end{concepto}

\column{0.45\textwidth}
\begin{aplicacion}
\textbf{Marketing:} segmentar clientes por comportamiento de compra para campañas dirigidas.

\medskip
\textbf{Medicina:} agrupar pacientes con perfiles clínicos similares para tratamientos personalizados.

\medskip
\textbf{Imágenes:} cuantización de color: reducir una imagen de millones de colores a $k$ colores representativos (compresión).

\medskip
\textbf{Astronomía:} clasificar galaxias por forma y brillo en catálogos de millones de objetos celestes.

\medskip
\textbf{Transporte:} ubicar depósitos de bicicletas públicas para minimizar la distancia promedio de los usuarios.
\end{aplicacion}
\end{columns}
\end{frame}

%---------------------------------------------------------------
\begin{frame}{K-means: ¿cuántos clusters?}
\begin{columns}[T]
\column{0.48\textwidth}
\textbf{Método del codo:}

Graficar inercia $W$ vs.\ $k$. Buscar el punto donde la reducción se aplana.

\medskip
\[
k=1 \to k=2: \text{ gran caída}
\]
\[
k=5 \to k=6: \text{ caída mínima} \Rightarrow k=5
\]

\medskip
\textbf{Coeficiente de silueta} para cada punto $i$:
\[
s_i = \frac{b_i - a_i}{\max(a_i, b_i)} \in [-1, 1]
\]
$a_i$: distancia media a su cluster.\\
$b_i$: distancia media mínima a otro cluster.

\medskip
$\bar{s}$ cerca de 1: clusters bien separados.

\column{0.48\textwidth}
\begin{aplicacion}
\textbf{Retail:} un supermercado analiza tickets de compra. Con $k=3$ emerge: ``compradores semanales de grandes volúmenes'', ``compradores diarios de poco'', ``compradores de ocasión''. La curva del codo confirma $k=3$.

\medskip
\textbf{Recursos humanos:} agrupar empleados por desempeño y antigüedad. La silueta indica si los grupos son realmente distintos o si se solapan.
\end{aplicacion}

\medskip
\begin{atencion}
K-means asume clusters \textbf{esféricos} y de \textbf{tamaño similar}. Si los datos tienen formas irregulares o muy distintos tamaños, K-means falla. Ver DBSCAN o GMM.
\end{atencion}
\end{columns}
\end{frame}

%---------------------------------------------------------------
\begin{frame}{Clustering jerárquico}
\begin{columns}[T]
\column{0.52\textwidth}
\textbf{Idea:} construir una jerarquía de fusiones.

\medskip
\textbf{Algoritmo aglomerativo:}
\begin{enumerate}
  \item Cada punto es su propio cluster.
  \item Fusionar los dos clusters más cercanos.
  \item Repetir hasta un único cluster.
\end{enumerate}

\medskip
\textbf{Métodos de enlace:}

{\small
\begin{tabular}{ll}
\toprule
\textbf{Enlace} & \textbf{Distancia entre $A$ y $B$} \\
\midrule
Simple   & mínima entre pares \\
Completo & máxima entre pares \\
Promedio & media de todas las distancias \\
Ward     & minimiza aumento de inercia \\
\bottomrule
\end{tabular}
}

\medskip
\begin{concepto}
El resultado es un \textbf{dendrograma}. El número de clusters $k$ se elige \emph{después} cortando el árbol a la altura adecuada.
\end{concepto}

\column{0.45\textwidth}
\begin{aplicacion}
\textbf{Genómica:} agrupar genes con patrones de expresión similares en experimentos de microarray. El dendrograma revela familias de genes que se activan juntos.

\medskip
\textbf{Lingüística:} construir familias de lenguas por similitud léxica. El árbol resultante coincide con clasificaciones históricas (romance, germánicas, eslavas).

\medskip
\textbf{Arqueología:} clasificar artefactos por composición química para inferir rutas de comercio.

\medskip
\textbf{Ecología:} agrupar especies por nicho ecológico a partir de variables ambientales.

\medskip
\textbf{Finanzas:} agrupar activos financieros para construir carteras diversificadas: clusters de acciones que se mueven juntas.
\end{aplicacion}
\end{columns}
\end{frame}

%---------------------------------------------------------------
\begin{frame}{DBSCAN: clustering por densidad}
\begin{columns}[T]
\column{0.52\textwidth}
\textbf{Parámetros:} radio $\varepsilon$ y mínimo de vecinos $m$.

\medskip
\textbf{Tipos de puntos:}
\begin{itemize}
  \item \textbf{Núcleo:} tiene $\geq m$ vecinos en radio $\varepsilon$.
  \item \textbf{Borde:} está en la vecindad de un núcleo.
  \item \textbf{Ruido:} ni núcleo ni borde.
\end{itemize}

\medskip
Los puntos núcleo conectados forman un cluster.

\medskip
\begin{concepto}
DBSCAN no requiere especificar $k$ y puede encontrar clusters de \textbf{cualquier forma}. Los outliers quedan etiquetados explícitamente como ruido.
\end{concepto}

\medskip
\textbf{Desventaja:} dificultoso con clusters de densidades muy distintas y con alta dimensionalidad.

\column{0.45\textwidth}
\begin{aplicacion}
\textbf{Geolocalización:} detectar zonas de alta concentración de accidentes de tráfico en una ciudad. La forma irregular de las zonas peligrosas no es capturada por K-means.

\medskip
\textbf{Astronomía:} identificar cúmulos de estrellas en mapas del cielo. Las galaxias tienen formas filamentosas, no esféricas.

\medskip
\textbf{Detección de fraude:} transacciones bancarias inusuales aparecen como ruido (outliers) rodeadas de clusters de comportamiento normal.

\medskip
\textbf{Manufactura:} detectar zonas defectuosas en superficies industriales a partir de mediciones de sensores.

\medskip
\textbf{Epidemiología:} identificar brotes geográficos de enfermedades infecciosas.
\end{aplicacion}
\end{columns}
\end{frame}

%---------------------------------------------------------------
\begin{frame}{Modelos de Mezclas Gaussianas (GMM)}
\begin{columns}[T]
\column{0.52\textwidth}
\textbf{Modelo:} los datos provienen de $k$ gaussianas:
\[
p(\mathbf{x}) = \sum_{l=1}^{k} \pi_l \,\mathcal{N}(\mathbf{x} \mid \boldsymbol{\mu}_l, \boldsymbol{\Sigma}_l)
\]

\medskip
\textbf{Estimación por EM:}
\begin{enumerate}
  \item \textbf{Paso E:} calcular probabilidad de pertenencia $r_{il}$ a cada componente.
  \item \textbf{Paso M:} actualizar $\pi_l$, $\boldsymbol{\mu}_l$, $\boldsymbol{\Sigma}_l$.
  \item Repetir hasta convergencia.
\end{enumerate}

\medskip
\begin{concepto}
A diferencia de K-means, la asignación es \textbf{blanda}: cada punto tiene probabilidad $r_{il}$ de pertenecer a cada cluster. K-means es un caso límite de GMM.
\end{concepto}

\medskip
Selección de $k$: BIC o AIC.

\column{0.45\textwidth}
\begin{aplicacion}
\textbf{Reconocimiento de voz:} modelar fonemas como mezclas de gaussianas en el espacio de características acústicas.

\medskip
\textbf{Biometría:} modelar distribuciones de huellas dactilares o iris para sistemas de identificación.

\medskip
\textbf{Finanzas:} modelar retornos financieros con régimen mixto: un componente de ``días normales'' y otro de ``crisis'' (alta volatilidad).

\medskip
\textbf{Medicina:} en imágenes de resonancia magnética, separar tejidos (materia gris, blanca, líquido cefalorraquídeo) cuyas intensidades de pixel siguen distribuciones gaussianas distintas.

\medskip
\textbf{Robótica:} modelar distribuciones de posición de objetos para robots que manipulan objetos en entornos inciertos.
\end{aplicacion}
\end{columns}
\end{frame}

%---------------------------------------------------------------
\begin{frame}{Evaluación de clustering}
Sin etiquetas verdaderas, la evaluación es \textbf{intrínseca}.

\medskip
\begin{columns}[T]
\column{0.48\textwidth}
\textbf{Coeficiente de silueta} $\bar{s} \in [-1,1]$

Para cada punto $i$:
\[
s_i = \frac{b_i - a_i}{\max(a_i,b_i)}
\]
$a_i$: distancia media dentro del cluster.\\
$b_i$: distancia media mínima al cluster más cercano.

\medskip
\textbf{Davies-Bouldin} (menor es mejor):
\[
DB = \frac{1}{k}\sum_{l}\max_{j\neq l}\frac{s_l + s_j}{d(\boldsymbol{\mu}_l,\boldsymbol{\mu}_j)}
\]

\textbf{Calinski-Harabasz} (mayor es mejor):
\[
CH = \frac{\text{traza}(B_k)/(k-1)}{\text{traza}(W_k)/(n-k)}
\]

\column{0.48\textwidth}
\begin{concepto}
No existe una única métrica ``correcta''. Lo ideal es usar varias y combinarlas con el criterio del dominio.
\end{concepto}

\medskip
\begin{aplicacion}
\textbf{Salud:} al segmentar pacientes, la silueta indica si los grupos son clínicamente distinguibles. Si $\bar{s} < 0.3$, los grupos se solapan y quizás el número de clusters es incorrecto.

\medskip
\textbf{Marketing:} al segmentar clientes, se usa Calinski-Harabasz para comparar soluciones con $k=3$, 4 o 5 y elegir la más compacta y separada. Luego un experto del negocio valida si los segmentos ``tienen sentido''.
\end{aplicacion}

\medskip
\begin{atencion}
La evaluación cuantitativa siempre debe complementarse con \textbf{validación de dominio}: que los grupos sean interpretables y accionables.
\end{atencion}
\end{columns}
\end{frame}

%---------------------------------------------------------------
\begin{frame}{Resumen: Bloque 1 -- Clustering}
\begin{center}
\renewcommand{\arraystretch}{1.4}
{\small
\begin{tabular}{lllll}
\toprule
\textbf{Método} & \textbf{Criterio} & \textbf{Forma} & \textbf{Outliers} & \textbf{Cuándo usarlo} \\
\midrule
K-means    & Inercia mínima    & Esférica   & Sensible & Clusters bien separados, $k$ conocido \\
Jerárquico & Enlace entre grupos & Flexible & Sensible & Quiero explorar la jerarquía \\
DBSCAN     & Densidad local    & Arbitraria & Detecta  & Forma irregular, outliers presentes \\
GMM        & Verosimilitud     & Elíptica   & Moderado & Necesito probabilidades de pertenencia \\
\bottomrule
\end{tabular}
}
\end{center}

\medskip
\begin{columns}[T]
\column{0.48\textwidth}
\textbf{Regla práctica:}
\begin{itemize}
  \item Empezar siempre con K-means para tener una línea de base.
  \item Si los datos tienen ruido o forma irregular $\to$ DBSCAN.
  \item Si quiero jerarquía o no sé $k$ $\to$ jerárquico.
  \item Si necesito incertidumbre en la asignación $\to$ GMM.
\end{itemize}

\column{0.48\textwidth}
\textbf{Pipeline típico:}
\begin{enumerate}
  \item Escalar variables.
  \item Explorar con K-means + codo + silueta.
  \item Refinar con DBSCAN o jerárquico si es necesario.
  \item Evaluar con múltiples métricas + criterio de dominio.
\end{enumerate}
\end{columns}
\end{frame}

\end{document}
